\section*{Bab \ref{ch:b2:atomic-orbitals} --- Orbital Atom: Bentuk, Energi, dan Ukuran}

\begin{solution}{exo:b2:atomic-orbitals:1}
3p: $n = 3$, $l = 1$; satu \omterm{def:b2:atomic-orbitals:node}{simpul radial} ($n - l - 1$) dan satu \omterm{def:b2:atomic-orbitals:node}{simpul
sudut}. 4d: $n = 4$, $l = 2$; satu \omterm{def:b2:atomic-orbitals:node}{simpul radial} dan dua \omterm{def:b2:atomic-orbitals:node}{simpul sudut}.
\end{solution}

\begin{solution}{exo:b2:atomic-orbitals:2}
$P_{2p} \propto r^4\mathrm e^{-r}$ ($r$ dalam $a_0$): $d\ln P/dr = 4/r - 1 = 0$
untuk $r = 4$: $4a_0 = \qty{212}{pm}$.
\end{solution}

\begin{solution}{exo:b2:atomic-orbitals:3}
$1 - \mathrm e^{-2}(1 + 2 + 2) = 1 - 5\mathrm e^{-2} = 0.323$.
\end{solution}

\begin{solution}{exo:b2:atomic-orbitals:4}
Empat cuping yang mengarah di antara sumbu $x$ dan $y$, pada bidang $xy$,
dengan tanda berselang-seling (positif di tempat $xy > 0$); bidang-bidang
simpulnya adalah $xz$ dan $yz$. \omterm{def:b2:atomic-orbitals:node}{Simpul radial}: $n - l - 1 = 0$.
\end{solution}

\begin{solution}{exo:b2:atomic-orbitals:5}
Oksigen, $(1\mathrm s)^2(2\mathrm s, 2\mathrm p)^6$: $\sigma = 5 \times 0.35 + 2 \times 0.85
= 3.45$, $Z^* = 4.55$; jari-jari $4/4.55 = 0.88a_0$, \qty{47}{pm}.
\end{solution}

\begin{solution}{exo:b2:atomic-orbitals:6}
Ionisasi mengambil elektron yang paling lemah terikat. Elektron-elektron
4s (\qty{-14}{eV}) jauh lebih lemah terikat daripada elektron-elektron 3d
(\qty{-59}{eV}): kedua elektron 4s pergi lebih dahulu, lalu satu elektron
3d, sehingga diperoleh $[\ce{Ar}]\,3\mathrm d^5$.
\end{solution}

\begin{solution}{exo:b2:atomic-orbitals:7}
Litium, 2s: $Z^* = 3 - 2 \times 0.85 = 1.30$, $I = 13.6 \times 1.30^2/4 =
\qty{5.75}{eV}$ (terukur 5.39, 7~\% terlalu tinggi). Natrium, 3s:
$Z^* = 11 - 8 \times 0.85 -
2 \times 1.00 = 2.20$,
$I = 13.6 \times 2.20^2/9 = \qty{7.3}{eV}$ (terukur 5.14): aturan-aturan
itu, sebuah model kasar, melebih-lebihkan ikatan elektron terluar, makin
lama makin besar menuruni golongan.
\end{solution}

\begin{solution}{exo:b2:atomic-orbitals:8}
Natrium: $Z^* = 2.20$, jari-jari $9/2.20 = 4.1a_0$ (\qty{216}{pm}). Klor,
3p: $\sigma = 6 \times 0.35 + 8 \times 0.85 + 2 \times 1.00 = 10.90$,
$Z^* = 6.10$, jari-jari $9/6.10 = 1.5a_0$ (\qty{78}{pm}). Kulitnya sama,
tetapi elektron terluar klor merasakan muatan hampir tiga kali lebih besar,
dan tertarik ke dalam.
\end{solution}

\begin{solution}{exo:b2:atomic-orbitals:9}
Keduanya mempunyai \omterm{def:b2:atomic-orbitals:radial-angular}{bagian sudut} $1/\sqrt{4\pi}$, sehingga integralnya menjadi
\[
  \int_0^\infty R_{1s}R_{2s}r^2\,dr \propto \int_0^\infty (2 - r)r^2\mathrm e^{-3r/2}\,dr =
  2 \times \frac{2!}{(3/2)^3} - \frac{3!}{(3/2)^4} = \frac{32}{27} - \frac{32}{27} = 0 .
\]
\end{solution}

\begin{solution}{exo:b2:atomic-orbitals:10}
$\int_0^\infty N^2r^2\mathrm e^{-r}r^2\,dr = N^2 \times 4! = 1$: $N = 1/\sqrt{24} =
1/(2\sqrt6)$, seperti pada tabel.
\end{solution}

\begin{solution}{exo:b2:atomic-orbitals:11}
$\langle r\rangle = 1.5a_0$, \omterm{def:b2:atomic-orbitals:radial-distribution}{jari-jari paling mungkin} $a_0$.
$P(r) = 4r^2\mathrm e^{-2r}$ naik tajam dari nol dan turun perlahan: ekor
panjang pada $r$ besar lebih berbobot dalam rata-rata daripada dalam posisi
maksimum.
\end{solution}

\begin{solution}{exo:b2:atomic-orbitals:12}
$Y_{p_z} \propto z/r = \cos\theta$ lenyap untuk $\theta = \pi/2$: bidang $xy$.
$\psi$ bertanda sama dengan $z$, positif di atas, negatif di bawah. Pada
bidang itu \omterm{def:b2:atomic-orbitals:wavefunction}{kepadatan peluangnya} nol.
\end{solution}

\begin{solution}{pb:b2:atomic-orbitals:1}
\textbf{1.} $\int|\psi|^2\,dV = \frac{4\pi}{\pi a_0^3}\int_0^\infty r^2\mathrm e^{-2r/a_0}\,dr =
\frac{4}{a_0^3} \times \frac{a_0^3}{4} = 1$.
\textbf{2.} $P(r) = 4\pi r^2|\psi|^2 = (4/a_0^3)\,r^2\mathrm e^{-2r/a_0}$.
\textbf{3.} $dP/dr = 0$: $2/r - 2/a_0 = 0$, $r = a_0$.
\textbf{4.} $\langle r\rangle = (4/a_0^3)\int_0^\infty r^3\mathrm e^{-2r/a_0}\,dr = (4/a_0^3)
\times 6(a_0/2)^4 = 1.5a_0$.
\textbf{5.} Distribusinya miring, dengan ekor panjang pada $r$ besar.
\textbf{6.} $|\psi|^2$ adalah kepadatan per satuan volume, yang paling besar
di inti; $P(r)$ memperhitungkan volume kulit, $4\pi r^2\,dr$, yang lenyap
pada $r = 0$.
\textbf{7.} \qty{52.9}{pm}.
\textbf{8.} Dengan $x = r/a_0$: $\int_{x_0}^\infty 4x^2\mathrm e^{-2x}\,dx$;
dengan integrasi parsial dua kali, $= \mathrm e^{-2x_0}(2x_0^2 + 2x_0 + 1)$.
\textbf{9.} $5\mathrm e^{-2} = 0.677$.
\textbf{10.} $13\mathrm e^{-4} = 0.238$.
\textbf{11.} $\mathrm e^{-2x}(1 + 2x + 2x^2) = 0.10$ untuk $x \approx 2.66$:
\qty{141}{pm}.
\textbf{12.} Kepadatannya tidak pernah lenyap pada jarak berapa pun: tidak
ada tepi. Tetapi 90~\% peluangnya terletak di dalam \qty{141}{pm}: sebuah
ukuran praktis.
\textbf{13.} 2s: satu \omterm{def:b2:atomic-orbitals:node}{simpul radial} (bola berjari-jari $2a_0$), tanpa \omterm{def:b2:atomic-orbitals:node}{simpul
sudut}. 2p: satu \omterm{def:b2:atomic-orbitals:node}{simpul sudut} (sebuah bidang), tanpa \omterm{def:b2:atomic-orbitals:node}{simpul radial}.
\textbf{14.} Pada $r = 2a_0$.
\textbf{15.} $4a_0$ (latihan 2).
\textbf{16.} $d\ln P/dr = 2/r - 2/(2 - r) - 1 = 0$ memberikan $r^2 - 6r + 4 = 0$,
$r = 3 \pm \sqrt5$: $0.76a_0$ (maksimum dalam yang kecil) dan $5.24a_0$
(maksimum utama).
\textbf{17.} 2s, melalui maksimum dalamnya. Dalam atom berelektron banyak,
elektron 2s menembus kulit 1s, kurang ditapis daripada elektron 2p, dan
energinya lebih rendah.
\textbf{18.} Semua jari-jari dibagi dengan $Z = 6$: 2p pada $0.67a_0$
(\qty{35}{pm}); maksimum 2s pada $0.13a_0$ dan $0.87a_0$.
\textbf{19.} C 3.25, N 3.90, O 4.55.
\textbf{20.} $\varepsilon = -13.6\,Z^{*2}/4$: C \qty{-35.9}{eV}, N \qty{-51.7}{eV}, O
\qty{-70.4}{eV}.
\textbf{21.} Seperti dalam bab ini: \qty{11.5}{eV} (terukur 11.26).
\textbf{22.} N: atom $5 \times (-13.6 \times 3.90^2/4) = \qty{-258.6}{eV}$; \ce{N+}, $Z^* =
4.25$, $4 \times (-13.6 \times 4.25^2/4) = \qty{-245.7}{eV}$; $I = \qty{12.9}{eV}$.
\textbf{23.} O: atom $6 \times (-13.6 \times 4.55^2/4) = \qty{-422.3}{eV}$; \ce{O+}, $Z^* =
4.90$, $5 \times (-13.6 \times 4.90^2/4) = \qty{-408.2}{eV}$; $I = \qty{14.2}{eV}$.
Model 11.5, 12.9, 14.2; terukur 11.26, 14.53, 13.62.
\textbf{24.} Energi terukur oksigen lebih rendah daripada energi nitrogen:
dalam oksigen, elektron 2p keempat harus berbagi orbital dengan elektron
lain, dan tolakan keduanya membuatnya lebih mudah dilepaskan, sedangkan
ketiga elektron p nitrogen yang tidak berpasangan membentuk subkulit
setengah penuh yang stabil. \omterm{def:b2:atomic-orbitals:slater}{Aturan Slater} memperlakukan semua elektron
dalam satu kelompok secara sama dan tidak mengenal pemasangan.
\textbf{25.} Elektron 1s hidrogen berada lebih jauh daripada $2a_0$ dengan
peluang $\boldsymbol{13\mathrm e^{-4} \approx 0.238}$.
\end{solution}
